% !TeX program = lualatex
\documentclass[11pt,a4paper]{article}
\usepackage[ngerman]{babel}
\usepackage{fontspec}
\usepackage[top=1.5cm,bottom=1.5cm,left=2cm,right=2cm]{geometry}
\usepackage{tcolorbox}
\tcbuselibrary{skins,breakable}
\usepackage{enumitem}
\usepackage{array}
\usepackage{booktabs}
\usepackage{multirow}
\usepackage{tikz}
\usetikzlibrary{arrows.meta,positioning,decorations.pathmorphing}
\usepackage{siunitx}
\sisetup{locale=DE,per-mode=symbol}
\usepackage{graphicx}
\usepackage{lastpage}
\usepackage{qrcode}
\usepackage[hidelinks]{hyperref}

% Farben
\definecolor{physikblau}{RGB}{44,90,160}
\definecolor{waermerot}{RGB}{211,47,47}
\definecolor{merkgruen}{RGB}{42,122,75}
\definecolor{hellgrau}{RGB}{245,245,245}

% Box-Definitionen (druckfreundlich)
\newtcolorbox{aufgabenbox}[1][]{
    colback=white,
    colframe=waermerot,
    coltitle=black,
    colbacktitle=white,
    fonttitle=\bfseries,
    boxrule=1pt,
    arc=0pt,
    left=8pt, right=8pt, top=6pt, bottom=6pt,
    title={#1}
}

\newtcolorbox{merkbox}[1][Merke]{
    colback=white,
    colframe=merkgruen,
    coltitle=black,
    colbacktitle=white,
    fonttitle=\bfseries\small,
    boxrule=1pt,
    arc=0pt,
    left=8pt, right=8pt, top=6pt, bottom=6pt,
    title={#1}
}

\newtcolorbox{protokollbox}[1][]{
    colback=white,
    colframe=physikblau,
    boxrule=0.5pt,
    arc=0pt,
    left=5pt, right=5pt, top=5pt, bottom=5pt,
    #1
}

% Lücken
\newcommand{\luecke}[1][3cm]{\underline{\hspace{#1}}}
\newcommand{\lueckelang}[1][6cm]{\underline{\hspace{#1}}}

% Niveau-Sterne
\newcommand{\niveauI}{$\star$}
\newcommand{\niveauII}{$\star\star$}
\newcommand{\niveauIII}{$\star\star\star$}

% Kopfzeile
\usepackage{fancyhdr}
\pagestyle{fancy}
\fancyhf{}
\fancyhead[L]{\small Physik Klasse 8}
\fancyhead[C]{\small Wärmelehre: Grundlagen}
\fancyhead[R]{\small AB-05}
\fancyfoot[C]{\small Seite \thepage\ von \pageref{LastPage}}
\renewcommand{\headrulewidth}{0.4pt}

\setlength{\parindent}{0pt}
\setlength{\parskip}{6pt}

\begin{document}

% Titelbereich
\noindent
\begin{minipage}[t]{0.82\textwidth}
    \vspace{0pt}
    {\Large\bfseries Arbeitsblatt: Wärmelehre -- Grundlagen}\\[3pt]
    {\normalsize Teilchenmodell $\bullet$ Aggregatzustände $\bullet$ Temperatur $\bullet$ Phasenübergänge}\\[8pt]
    \begin{tabular}{|p{4cm}|p{4cm}|p{3cm}|}
        \hline
        Name: & Klasse: & Datum: \\[8pt]
        \hline
    \end{tabular}
\end{minipage}
\hfill
\begin{minipage}[t]{0.15\textwidth}
    \vspace{0pt}
    \raggedleft
    \qrcode[height=1.4cm]{https://devbox42.github.io/sciencesim/physik/kl08/05-waerme/LP-05-waerme-grundlagen.html}%
    \llap{\href{https://devbox42.github.io/sciencesim/physik/kl08/05-waerme/LP-05-waerme-grundlagen.html}{\phantom{\rule{1.4cm}{1.4cm}}}}
\end{minipage}

\vspace{5pt}

% ===== AUFGABE 1 =====
\begin{aufgabenbox}[Aufgabe 1: Einstieg -- Wärme im Alltag \hfill \niveauI]
Notiere drei Situationen aus deinem Alltag, in denen etwas warm oder kalt wird.

\begin{enumerate}[label=\alph*)]
    \item \lueckelang[10cm]
    \item \lueckelang[10cm]
    \item \lueckelang[10cm]
\end{enumerate}

Was haben alle diese Situationen gemeinsam?

\lueckelang[12cm]
\end{aufgabenbox}

% ===== AUFGABE 2 =====
\begin{aufgabenbox}[Aufgabe 2: Das Teilchenmodell \hfill \niveauI]
Vervollständige die drei Grundaussagen des Teilchenmodells:

\begin{enumerate}[label=\arabic*.]
    \item Alle Stoffe bestehen aus \luecke[4cm].
    \item Diese Teilchen sind in \luecke[4cm] Bewegung.
    \item Zwischen den Teilchen gibt es \luecke[4cm].
\end{enumerate}
\end{aufgabenbox}

\begin{merkbox}[Merke: Teilchenmodell]
Die Bewegung der Teilchen nennt man \textbf{thermische Bewegung}. Sie hört nie auf -- selbst wenn ein Gegenstand völlig ruhig erscheint, bewegen sich die Teilchen darin ständig!
\end{merkbox}

% ===== AUFGABE 3 =====
\begin{aufgabenbox}[Aufgabe 3: Die drei Aggregatzustände \hfill \niveauI/\niveauII]
Zeichne die Teilchenanordnung für jeden Aggregatzustand und beschreibe die Bewegung.

\begin{center}
\begin{tabular}{|c|c|c|}
\hline
\textbf{FEST} & \textbf{FLÜSSIG} & \textbf{GASFÖRMIG} \\
\hline
\begin{tikzpicture}
\draw[thick] (0,0) rectangle (3,3);
\node at (1.5,1.5) {\small Zeichnung};
\end{tikzpicture}
&
\begin{tikzpicture}
\draw[thick] (0,0) rectangle (3,3);
\node at (1.5,1.5) {\small Zeichnung};
\end{tikzpicture}
&
\begin{tikzpicture}
\draw[thick] (0,0) rectangle (3,3);
\node at (1.5,1.5) {\small Zeichnung};
\end{tikzpicture}
\\
\hline
\begin{minipage}{3cm}\centering\small
Teilchen\\[3pt]
\luecke[2.5cm]\\[3pt]
um feste Plätze
\end{minipage}
&
\begin{minipage}{3cm}\centering\small
Teilchen\\[3pt]
\luecke[2.5cm]\\[3pt]
aneinander vorbei
\end{minipage}
&
\begin{minipage}{3cm}\centering\small
Teilchen\\[3pt]
\luecke[2.5cm]\\[3pt]
frei umher
\end{minipage}
\\
\hline
\end{tabular}
\end{center}

\textbf{Bindung:} Je \luecke[3cm] die Bindungskräfte zwischen den Teilchen, desto \luecke[3cm] ist der Stoff.
\end{aufgabenbox}

% ===== AUFGABE 4 =====
\begin{aufgabenbox}[Aufgabe 4: Was ist Temperatur? \hfill \niveauII]
Erkläre in eigenen Worten, was Temperatur auf Teilchenebene bedeutet.

\begin{protokollbox}
\vspace{2cm}
\end{protokollbox}

Vervollständige:
\begin{itemize}[leftmargin=*]
    \item Hohe Temperatur $\rightarrow$ Teilchen bewegen sich \luecke[3cm] $\rightarrow$ \luecke[2cm] kinetische Energie
    \item Niedrige Temperatur $\rightarrow$ Teilchen bewegen sich \luecke[3cm] $\rightarrow$ \luecke[2cm] kinetische Energie
\end{itemize}
\end{aufgabenbox}

\begin{merkbox}[Merke: Temperatur]
\textbf{Temperatur} ist ein Maß für die \lueckelang[4cm] der Teilchen eines Stoffes.
\end{merkbox}

% ===== AUFGABE 5 =====
\begin{aufgabenbox}[Aufgabe 5: Temperaturskalen -- Celsius und Kelvin \hfill \niveauI/\niveauII]

\textbf{a) Fixpunkte der Celsius-Skala:}
\begin{itemize}[leftmargin=*]
    \item \SI{0}{\degreeCelsius} = \lueckelang[5cm]
    \item \SI{100}{\degreeCelsius} = \lueckelang[5cm]
\end{itemize}

\textbf{b) Die Kelvin-Skala:}

Der \textbf{absolute Nullpunkt} liegt bei \luecke[2cm] K = \luecke[2cm] °C.

Bei dieser Temperatur bewegen sich die Teilchen (fast) \luecke[3cm].

\textbf{c) Umrechnung:}

\begin{center}
\fbox{\Large $T$ / K = $\vartheta$ / °C + 273}
\end{center}

\textbf{d) Rechne um:}
\begin{enumerate}[label=\roman*)]
    \item \SI{25}{\degreeCelsius} = \luecke[2cm] K
    \item \SI{350}{K} = \luecke[2cm] °C
    \item \SI{-40}{\degreeCelsius} = \luecke[2cm] K
    \item \SI{0}{K} = \luecke[2cm] °C
\end{enumerate}
\end{aufgabenbox}

% ===== AUFGABE 6 =====
\begin{aufgabenbox}[Aufgabe 6: Erweiterung -- Fahrenheit \hfill \niveauIII]
In den USA wird die Fahrenheit-Skala verwendet.

\smallskip
\begin{center}
\fbox{\;\; $\displaystyle \vartheta_F = \frac{9}{5} \cdot \vartheta_C + 32$ \qquad\qquad $\displaystyle \vartheta_C = \frac{5}{9} \cdot (\vartheta_F - 32)$ \;\;}
\end{center}
\smallskip

\textbf{a)} Wasser gefriert bei \luecke[2cm] °F und siedet bei \luecke[2cm] °F.

\textbf{b)} Bei welcher Temperatur sind Celsius und Fahrenheit gleich? \luecke[2cm]

\textbf{c)} Recherche: Warum hat sich Fahrenheit diese ungewöhnlichen Fixpunkte ausgedacht?

\begin{protokollbox}
\vspace{1.5cm}
\end{protokollbox}
\end{aufgabenbox}

% ===== AUFGABE 7 =====
\begin{aufgabenbox}[Aufgabe 7: Die sechs Phasenübergänge \hfill \niveauI/\niveauII]
Beschrifte die Pfeile mit den richtigen Begriffen:

\begin{center}
\begin{tikzpicture}[
    box/.style={draw, thick, minimum width=2.2cm, minimum height=1cm, align=center, rounded corners=4pt},
    arr/.style={-{Stealth[length=3mm]}, thick}
]
% Boxes (triangular: gas top, fest bottom-left, flüssig bottom-right)
\node[box, fill=orange!10] (gas) at (3.5,4) {GASFÖRMIG};
\node[box, fill=blue!10] (fest) at (0,0) {FEST};
\node[box, fill=green!10] (fluessig) at (7,0) {FLÜSSIG};

% Bottom: Schmelzen / Erstarren
\draw[arr, waermerot] ([yshift=3pt]fest.east) -- node[above, font=\small] {\luecke[1.8cm]} ([yshift=3pt]fluessig.west);
\draw[arr, physikblau] ([yshift=-3pt]fluessig.west) -- node[below, font=\small] {\luecke[1.8cm]} ([yshift=-3pt]fest.east);

% Left: Sublimieren / Resublimieren
\draw[arr, purple, bend left=8] (fest.north) to node[left, font=\small] {\luecke[1.8cm]} (gas.south west);
\draw[arr, teal, bend left=8] (gas.south west) to node[right, font=\small] {\luecke[1.8cm]} (fest.north);

% Right: Verdampfen / Kondensieren
\draw[arr, waermerot, bend right=8] (fluessig.north) to node[right, font=\small] {\luecke[1.8cm]} (gas.south east);
\draw[arr, physikblau, bend right=8] (gas.south east) to node[left, font=\small] {\luecke[1.8cm]} (fluessig.north);

% Legend
\node[font=\scriptsize, text=gray] at (3.5,-0.7) {rot = Energie zuführen \quad blau = Energie abgeben};
\end{tikzpicture}
\end{center}

\textbf{Begriffe:} Schmelzen, Erstarren, Verdampfen, Kondensieren, Sublimieren, Resublimieren

\textbf{Energie:} Bei welchen Übergängen wird Energie \textbf{zugeführt}? \lueckelang[6cm]

Bei welchen Übergängen wird Energie \textbf{abgegeben}? \lueckelang[6cm]
\end{aufgabenbox}

\newpage

% ===== AUFGABE 8 =====
\begin{aufgabenbox}[Aufgabe 8: Simulationsexperiment -- Erwärmung von Eis \hfill \niveauII]

\textbf{Simulation:} Im Lernpfad (Schritt 8) siehst du ein Becherglas mit Eis auf einer Heizplatte. Ein Thermometer zeigt die Temperatur, eine Uhr läuft mit.

\textbf{Auftrag:} Lies alle \SI{30}{s} die Temperatur ab und notiere den Messwert.

\renewcommand{\arraystretch}{1.15}
\begin{center}
\begin{tabular}{|c|p{1.5cm}||c|p{1.5cm}|}
\hline
\textbf{Zeit} & \textbf{$\vartheta$ / °C} & \textbf{Zeit} & \textbf{$\vartheta$ / °C} \\
\hline
0:00 & & 3:00 & \\
\hline
0:30 & & 3:30 & \\
\hline
1:00 & & 4:00 & \\
\hline
1:30 & & 4:30 & \\
\hline
2:00 & & 5:00 & \\
\hline
2:30 & & & \\
\hline
\end{tabular}
\end{center}

Beobachtungen während der Simulation: \lueckelang[9cm]

\vspace{4pt}
\textbf{Trage deine Messwerte als Punkte in das Diagramm ein und verbinde sie:}

\begin{center}
\begin{tikzpicture}
    % Light gray background
    \fill[gray!5] (0,-1) rectangle (10,6);
    % Fine grid (5mm squares)
    \draw[step=0.5, very thin, gray!25] (0,-1) grid (10,6);
    % Major grid (1cm)
    \draw[step=1, thin, gray!50] (0,-1) grid (10,6);
    % Reference lines at 0°C and 100°C
    \draw[dashed, blue!40, semithick] (0,0) -- (10,0);
    \draw[dashed, red!40, semithick] (0,5) -- (10,5);
    \node[right, font=\tiny, blue!60] at (10.1,0) {0\,°C};
    \node[right, font=\tiny, red!60] at (10.1,5) {100\,°C};
    % Y-axis
    \draw[very thick, -{Stealth[length=3mm]}] (0,-1) -- (0,6.5);
    \node[above, font=\small] at (0,6.5) {$\vartheta$ / °C};
    % X-axis
    \draw[very thick, -{Stealth[length=3mm]}] (0,-1) -- (10.5,-1);
    \node[below, font=\small] at (5,-1.6) {Zeit $t$ / min};
    % X labels (every minute)
    \foreach \x/\lab in {0/0, 2/1, 4/2, 6/3, 8/4, 10/5} {
        \draw[thick] (\x,-1) -- (\x,-1.15);
        \node[below, font=\scriptsize] at (\x,-1.15) {\lab};
    }
    % X minor ticks (every 30s)
    \foreach \x in {1,3,5,7,9} {
        \draw (\x,-1) -- (\x,-1.08);
    }
    % Y labels (every 20°C)
    \foreach \y/\temp in {-1/-20, 0/0, 1/20, 2/40, 3/60, 4/80, 5/100, 6/120} {
        \draw[thick] (0,\y) -- (-0.15,\y);
        \node[left, font=\scriptsize] at (-0.15,\y) {\temp};
    }
\end{tikzpicture}
\end{center}

\textbf{a)} Markiere den Bereich, in dem die Temperatur \textbf{konstant} bleibt (Plateau). Bei welcher Temperatur liegt das Plateau? \luecke[2cm]

\textbf{b)} Beschreibe den Verlauf deiner Kurve und erkläre, warum sie so aussieht:

\begin{protokollbox}
\vspace{1.5cm}
\end{protokollbox}

\textit{\small Tipp: Die Messpunkte in der Simulation helfen dir beim Ablesen.}
\end{aufgabenbox}

% ===== AUFGABE 9 =====
\begin{aufgabenbox}[Aufgabe 9: Verdunsten vs. Sieden \hfill \niveauII]
Erkläre den Unterschied zwischen Verdunsten und Sieden:

\begin{center}
\renewcommand{\arraystretch}{1.5}
\begin{tabular}{|p{3cm}|p{5cm}|p{5cm}|}
\hline
& \textbf{Verdunsten} & \textbf{Sieden} \\
\hline
\textbf{Wo?} & & \\
\hline
\textbf{Temperatur?} & & \\
\hline
\textbf{Geschwindigkeit?} & & \\
\hline
\textbf{Beispiel:} & & \\
\hline
\end{tabular}
\end{center}

\textbf{Anwendung:} Erkläre, warum nasse Haut bei Wind kühler wird als trockene Haut:

\begin{protokollbox}
\vspace{2cm}
\end{protokollbox}
\end{aufgabenbox}

\newpage

% ===== AUFGABE 10 =====
\begin{aufgabenbox}[Aufgabe 10: Zusammenfassung -- Lückentext \hfill \niveauI]
Fülle die Lücken aus:

Alle Stoffe bestehen aus \luecke[3cm], die sich \luecke[3cm] bewegen. Diese Bewegung nennt man \luecke[4cm] Bewegung.

Je nachdem, wie stark die Teilchen aneinander gebunden sind, unterscheiden wir drei \luecke[4cm]: fest, flüssig und gasförmig.

Die \luecke[3cm] ist ein Maß für die mittlere \luecke[4cm] der Teilchen. Je schneller sich die Teilchen bewegen, desto \luecke[3cm] ist die Temperatur.

Die \luecke[3cm]-Skala beginnt beim absoluten Nullpunkt (\SI{0}{K} = \SI{-273}{\degreeCelsius}). Bei dieser Temperatur bewegen sich die Teilchen (fast) gar nicht mehr.

Wenn ein Stoff seinen Aggregatzustand ändert, nennt man das \luecke[4cm]. Während dieses Vorgangs bleibt die Temperatur \luecke[3cm], obwohl weiter Energie zugeführt wird. Diese Energie wird zum \luecke[4cm] der Bindungen zwischen den Teilchen verwendet.

\textbf{Lösungswörter:} Teilchen, ständig, thermische, Aggregatzustände, Temperatur, kinetische Energie, höher, Kelvin, Phasenübergang, konstant, Aufbrechen
\end{aufgabenbox}

% ===== ZUSATZAUFGABE =====
\begin{aufgabenbox}[Zusatzaufgabe: Transferaufgabe \hfill \niveauIII]
Ein Eiswürfel wird aus dem Gefrierfach (\SI{-18}{\degreeCelsius}) genommen und in ein Glas mit Wasser (\SI{20}{\degreeCelsius}) gelegt.

\textbf{a)} Beschreibe, was auf Teilchenebene passiert, wenn der Eiswürfel schmilzt.

\begin{protokollbox}
\vspace{2cm}
\end{protokollbox}

\textbf{b)} Warum wird das Wasser im Glas kälter?

\begin{protokollbox}
\vspace{1.5cm}
\end{protokollbox}

\textbf{c)} Bei welcher Temperatur würde der Schmelzvorgang aufhören, wenn man unbegrenzt viel Eis hinzufügt? Begründe.

\begin{protokollbox}
\vspace{1.5cm}
\end{protokollbox}
\end{aufgabenbox}

\vfill

\begin{center}
\small\textit{Viel Erfolg!}
\end{center}

\end{document}
